% ============================================================
% Preamble: formatting compliance with the DIT Thesis Guideline
% (Deggendorf Institute of Technology, "Thesis Guideline",
%  Version 1.0, 27.06.2023)
%
% Requirements implemented here:
%  - DIN A4, single-sided printing
%  - Font size 11-12pt body, 14pt chapter headings
%  - Margins: top 2.5cm, bottom 2.5cm, left 3cm, right 2cm
%  - 1.5 line spacing
%  - Justified text with automatic hyphenation
%  - Consecutive page numbers in header/footer
%  - Section numbering 1, 1.1, 1.2, 2, 2.1, 2.1.1 ...
%  - Headings: not underlined/indented, bold, same font as body,
%    start with capital letter, fit on one line
%  - No university logo anywhere (guideline 4.1.1 / 3.5)
% ============================================================

% --- Page geometry -------------------------------------------------
\usepackage[a4paper,top=2.5cm,bottom=2.5cm,left=3cm,right=2cm]{geometry}

% --- Language / hyphenation -----------------------------------------
% Switch babel's main language to ngerman if you write in German;
% keep english if the thesis is written in English (check with your
% supervisor which language is expected).
\usepackage[english]{babel}
\usepackage[utf8]{inputenc}
\usepackage[T1]{fontenc}

% Latin Modern: a fully scalable (outline) replacement for the default
% Computer Modern fonts. Needed because titlesec below requests
% non-standard font sizes (e.g. \fontsize{14}{16} for chapter
% headings); without lmodern, LaTeX falls back to Metafont-generated
% bitmap fonts at those sizes, which are not scalable - and microtype's
% font expansion feature then fails with a fatal
% "auto expansion is only possible with scalable fonts" error.
\usepackage{lmodern}

% --- Font size (11-12pt body, set via documentclass option) --------
% Body font size is set as a \documentclass option in main.tex.

% --- Line spacing: 1.5 -----------------------------------------------
\usepackage{setspace}
\onehalfspacing

% --- Justification with automatic hyphenation ------------------------
% LaTeX justifies by default; microtype improves justification quality
% (fewer visible stretched/compressed lines) without changing the
% guideline-required "justification with automatic hyphenation" look.
\usepackage{microtype}

% This thesis's prose is dense with long, non-hyphenating \texttt{}
% identifiers (ROS topic/plugin/parameter names). Without headroom,
% TeX sometimes has no legal break short of pushing text past the
% right margin ("Overfull \hbox"), which would visibly violate the
% guideline's "illustrations and tables [text] are to remain within
% the page margins" requirement. emergencystretch gives the
% line-breaker extra last-resort stretch before it would otherwise
% overfull, fixing these without loosening ordinary paragraphs.
\emergencystretch=3em

% A handful of table cells hold long filename-like identifiers
% (e.g. libgazebo_ros_ackermann_drive.so) that are still wider than
% their column even with emergencystretch, since that only adds
% inter-word stretch and these are single unbreakable "words". seqsplit
% allows breaking such a token at any character as a last resort, kept
% to just those specific cells rather than used document-wide.
\usepackage{seqsplit}

% --- Headings: bold, same font as body text, not underlined ----------
\usepackage{titlesec}
\titleformat{\chapter}[hang]
  {\normalfont\fontsize{14}{16}\bfseries}{\thechapter}{1em}{}
\titleformat{\section}
  {\normalfont\fontsize{12}{14}\bfseries}{\thesection}{1em}{}
\titleformat{\subsection}
  {\normalfont\fontsize{11}{13}\bfseries}{\thesubsection}{1em}{}
\titleformat{\subsubsection}
  {\normalfont\fontsize{11}{13}\bfseries}{\thesubsubsection}{1em}{}
\titlespacing*{\chapter}{0pt}{0pt}{20pt}

% --- Section numbering depth (down to subsubsection, e.g. 2.1.1) -----
\setcounter{secnumdepth}{3}
\setcounter{tocdepth}{3}

% --- Footnotes: numbered consecutively throughout the document, not
% reset per chapter (mandatory per the Campus Cham "Specific
% Instructions" registration document) -----------------------------
\usepackage{chngcntr}
\counterwithout{footnote}{chapter}

% --- Header/footer: consecutive page numbers --------------------------
\usepackage{fancyhdr}
\pagestyle{fancy}
\fancyhf{}
\fancyfoot[C]{\thepage}
\renewcommand{\headrulewidth}{0pt}

% A plain style for chapter-opening pages (still numbered, no header rule)
\fancypagestyle{plain}{
  \fancyhf{}
  \fancyfoot[C]{\thepage}
  \renewcommand{\headrulewidth}{0pt}
}

% --- Figures and tables ------------------------------------------------
\usepackage{graphicx}
\usepackage{caption}
\usepackage{subcaption}
\captionsetup{font=small,labelfont=bf,justification=justified}
\usepackage{booktabs}   % clean table rules (no vertical/frame lines)
\usepackage{longtable}  % tables that span multiple pages (Results ch.)
\usepackage{array}

% --- Diagrams and plots -------------------------------------------------
% No photos/screenshots exist for this thesis; every figure is instead
% drawn directly from data and mechanisms already described in the
% prose (architecture block diagrams, a path-geometry diagram, and
% pgfplots charts built from the thesis's own result tables), so a
% reader never has to take a described mechanism or number on faith.
\usepackage{tikz}
\usetikzlibrary{arrows.meta,positioning,calc,fit,backgrounds}
\usepackage{pgfplots}
\pgfplotsset{compat=1.17}

% --- Formulas and physical quantities -----------------------------------
\usepackage{amsmath}
\usepackage{amssymb}
\usepackage{siunitx}    % correct unit typesetting: \SI{2.3}{\ampere} etc.
\sisetup{
  detect-all,
  output-decimal-marker = {.},   % dot as decimal separator (English)
  per-mode = symbol
}

% --- Code / configuration listings (launch files, YAML, Python) --------
\usepackage{listings}
\usepackage{xcolor}
\definecolor{codebg}{RGB}{245,245,245}

% Custom, minimal language definitions for every language= used in this
% thesis's lstlisting blocks (Python, C++, XML, bash), defined BEFORE
% any \begin{lstlisting} uses them. This deliberately overrides
% `listings`'s own bundled per-language keyword databases
% (lstlang1.sty/lstlang2.sty/lstlang3.sty) rather than relying on them -
% a known TeX Live 2026 regression in that bundled data triggers
% "Extra alignment tab has been changed to \cr" from inside
% lstlang1.sty on some language definitions. Defining the language
% ourselves means `listings` never loads the bundled file for these
% four names at all, sidestepping the bug regardless of exactly which
% bundled definition is at fault. Highlighting is intentionally simple
% (comments, strings, a short keyword list) - plenty for short
% illustrative snippets, and it will never break a compile this way.
\lstdefinelanguage{Python}{
  morekeywords={def,class,if,elif,else,for,while,return,import,from,as,
    try,except,finally,with,lambda,None,True,False,and,or,not,in,is,
    break,continue,pass,raise,yield,global,nonlocal,self,print},
  sensitive=true,
  morecomment=[l]{\#},
  morestring=[b]",
  morestring=[b]'
}
\lstdefinelanguage{XML}{
  morecomment=[s]{<!--}{-->},
  morestring=[b]",
  morestring=[b]',
  sensitive=true
}
\lstdefinelanguage{C++}{
  morekeywords={void,int,float,double,bool,char,const,static,class,
    struct,namespace,public,private,protected,return,if,else,for,
    while,do,switch,case,break,continue,new,delete,template,typename,
    using,auto,size_t,true,false,nullptr,enum},
  sensitive=true,
  morecomment=[l]{//},
  morecomment=[s]{/*}{*/},
  morestring=[b]"
}
\lstdefinelanguage{bash}{
  morekeywords={sudo,python3,ros2,run,launch,env,export,echo,cd},
  sensitive=true,
  morecomment=[l]{\#},
  morestring=[b]",
  morestring=[b]'
}

\lstset{
  basicstyle=\ttfamily\small,
  backgroundcolor=\color{codebg},
  breaklines=true,
  frame=single,
  framerule=0.3pt,
  columns=flexible,
  captionpos=b,
  numbers=left,
  numberstyle=\tiny,
  keywordstyle=\color{blue!60!black}\bfseries,
  commentstyle=\color{gray},
  showstringspaces=false
}

% --- Bibliography: numeric, DIN-style [1]-style in-text citation -------
% Plain LaTeX + BibTeX (\bibliographystyle{plain} in main.tex), not
% biblatex/biber: functionally equivalent numeric-bracket output for a
% bibliography this size, but only one extra compile pass instead of
% biber's own (slower) separate backend - matters on Overleaf's free
% compile-time limit. \bibliographystyle{plain} already sorts entries
% alphabetically by author while numbering them [1], [2], ... in that
% order, matching the guideline's "numbered upon citation order OR
% alphabetically sorted" allowance (here: both at once).

% --- Abbreviations list --------------------------------------------------
% Rendered as a plain table (frontmatter/abbreviations.tex) rather than
% via the glossaries package, which needs its own external
% makeglossaries run - another full extra pass this document doesn't
% need for a short, static, alphabetically-sorted list.

% --- Hyperlinks (internal cross-references only; keep unobtrusive) -----
\usepackage[hidelinks]{hyperref}
\usepackage{cleveref}

% --- Appendix numbering: A1, A1.1, A1.2, A2, A3 ... ---------------------
\usepackage{appendix}
